\par\Needspace{6\baselineskip}\medskip
\hypertarget{native-137-article099}{}\pdfbookmark[2]{99.}{article-099}
\begin{center}99.\end{center}
In HYPERBOLA quoque characteres $p$, $v$, $v'$, $f$, $F$, $r$, $r'$, $t$ in significatione eadem retinemus, pro semiaxi maiori $a$ autem, qui hic negativus est, scribemus $-\alpha$; excentricitatem $e$ perinde ut supra art. 21 etc. statuemus $=\dfrac1{\cos\psi}\mathpunct{\cdot}$ Quantitatem auxiliarem illic per $u$ expressam, statuemus pro loco primo $=\dfrac Cc$, pro secundo $=Cc$, unde facile concluditur, $c$ semper esse maiorem quam 1, \SourceLeaf{138}sed ceteris paribus eo minus differre ab 1, quo minus duo loci propositi ab invicem distent. Ex aequationibus in art. 21 evolutis huc transferimus forma paullulum mutata sextam et septimam
\[
\begin{aligned}
\mathrm{1)}\qquad&\cos\tfrac12v=\tfrac12\left(\BareRoot{\dfrac Cc}+\BareRoot{\dfrac cC}\right)\BareRoot{\dfrac{(e-1)\alpha}r}\\[.8em]
\mathrm{2)}\qquad&\sin\tfrac12v=\tfrac12\left(\BareRoot{\dfrac Cc}-\BareRoot{\dfrac cC}\right)\BareRoot{\dfrac{(e+1)\alpha}r}\\[.8em]
\mathrm{3)}\qquad&\cos\tfrac12v'=\tfrac12\left(\BareRoot{Cc}+\BareRoot{\dfrac1{Cc}}\right)\BareRoot{\dfrac{(e-1)\alpha}{r'}}\\[.8em]
\mathrm{4)}\qquad&\sin\tfrac12v'=\tfrac12\left(\BareRoot{Cc}-\BareRoot{\dfrac1{Cc}}\right)\BareRoot{\dfrac{(e+1)\alpha}{r'}}\mathpunct{\cdot}
\end{aligned}
\]
\noindent Hinc statim demanant sequentes:
\[
\begin{aligned}
\mathrm{5)}\qquad&\sin F=\tfrac12\alpha\left(C-\dfrac1C\right)\BareRoot{\dfrac{ee-1}{rr'}}\\[.8em]
\mathrm{6)}\qquad&\sin f=\tfrac12\alpha\left(c-\dfrac1c\right)\BareRoot{\dfrac{ee-1}{rr'}}\\[.8em]
\mathrm{7)}\qquad&\cos F=\left(e\left(c+\dfrac1c\right)-\left(C+\dfrac1C\right)\right)\dfrac\alpha{2\BareRoot{rr'}}\\[.8em]
\mathrm{8)}\qquad&\cos f=\left(e\left(C+\dfrac1C\right)-\left(c+\dfrac1c\right)\right)\dfrac\alpha{2\BareRoot{rr'}}\mathpunct{\cdot}
\end{aligned}
\]
\noindent Porro fit per aequ. X art. 21
\[
\begin{aligned}
\frac r\alpha&=\tfrac12e\left(\frac Cc+\frac cC\right)-1\\[.6em]
\frac{r'}\alpha&=\tfrac12e\left(Cc+\frac1{Cc}\right)-1,
\end{aligned}
\]
\noindent atque hinc
\[
\begin{aligned}
\mathrm{9)}\qquad&\frac{r'-r}\alpha=\tfrac12e\left(C-\frac1C\right)\left(c-\frac1c\right)\\[.6em]
\mathrm{10)}\qquad&\frac{r'+r}\alpha=\tfrac12e\left(C+\frac1C\right)\left(c+\frac1c\right)-2.
\end{aligned}
\]
\noindent Haec aequatio 10 cum 8 combinata praebet
\[
\mathrm{11)}\qquad\alpha=\frac{r'+r-\left(c+\dfrac1c\right)\cos f\,\mathpunct{.}\BareRoot{rr'}}{\tfrac12\left(c-\dfrac1c\right)^2}\mathpunct{\cdot}
\]
\noindent Statuendo itaque perinde ut in ellipsi $\dfrac{\BareRoot{\dfrac{r'}r}+\BareRoot{\dfrac r{r'}}}{2\cos f}=1+2l$, vel $=1-2L$, \SourceLeaf{139}prout $\cos f$ est positivus vel negativus, fit
\[
\begin{aligned}
\mathrm{12)}\qquad&\alpha=\frac{8\left(l-\tfrac14\left(\BareRoot c-\BareRoot{\dfrac1c}\right)^2\right)\cos f\,\mathpunct{.}\BareRoot{rr'}}{\left(c-\dfrac1c\right)^2}\\[1em]
\mathrm{12^*)}\qquad&\alpha=\frac{-8\left(L+\tfrac14\left(\BareRoot c-\BareRoot{\dfrac1c}\right)^2\right)\cos f\,\mathpunct{.}\BareRoot{rr'}}{\left(c-\dfrac1c\right)^2}\mathpunct{\cdot}
\end{aligned}
\]
\noindent Computus quantitatis $l$ vel $L$ hic perinde ut in ellipsi adiumento anguli auxiliaris $\omega$ instituetur. Denique fit ex aequatione XI art. 22 (accipiendo logarithmos hyperbolicos)
\[
\begin{aligned}
\frac{kt}{\alpha^{\frac32}}&=\tfrac12e\left(Cc-\frac1{Cc}-\frac Cc+\frac cC\right)-\log Cc+\log\frac Cc\\[.6em]
&=\tfrac12e\left(C+\frac1C\right)\left(c-\frac1c\right)-2\log c
\end{aligned}
\]
\noindent sive eliminata $C$ adiumento aequationis 8
\[
\frac{kt}{\alpha^{\frac32}}=\frac{\left(c-\dfrac1c\right)\cos f\,\mathpunct{.}\BareRoot{rr'}}\alpha\mathpunct{\raisebox{-1ex}{.}}+\tfrac12\left(cc-\frac1{cc}\right)-2\log c.
\]
\noindent In hac aequatione pro $\alpha$ substituimus valorem eius ex 12, $12^*$; dein characterem $m$ vel $M$ in eadem significatione, quam formulae 11, $11^*$ art. 88 assignant, introducimus; tandemque brevitatis gratia scribimus
\[
\begin{gathered}
\tfrac14\left(\BareRoot c-\BareRoot{\frac1c}\right)^2=z,\\[.6em]
\frac{cc-\dfrac1{cc}-4\log c}{\tfrac14\left(c-\dfrac1c\right)^3}=Z,
\end{gathered}
\]
\noindent quo facto oriuntur aequationes
\[
\begin{aligned}
\mathrm{13)}\qquad&m=(l-z)^{\frac12}+(l-z)^{\frac32}Z\\[.6em]
\mathrm{13^*)}\qquad&M=-(L+z)^{\frac12}+(L+z)^{\frac32}Z,
\end{aligned}
\]
\noindent quae unicam incognitam $z$ implicant, quum manifesto sit $Z$ functio ipsius $z$ per formulam sequentem expressa
\[
Z=\frac{(1+2z)\BareRoot{(z+zz)}-\log(\BareRoot{1+z}+\BareRoot z)}{2(z+zz)^{\frac32}}\mathpunct{\cdot}
\]
